% ══════════════════════════════════════════════════════════════════════════
%  Chapter 2: The Boolean Hypercube and Fourier Analysis
% ══════════════════════════════════════════════════════════════════════════

\chapter{Fourier Analysis on the Boolean Hypercube}\label{ch:boolean-fourier}

This chapter develops the Fourier analysis of functions on $\pmone^n$
from scratch. The treatment follows \citet{odonnell2014analysis},
Chapters~1--2, which is the standard reference. A shorter introduction
is \citet{dewolf2008brief}.

% ──────────────────────────────────────────────────────────────────────
\section{The Boolean Hypercube}\label{sec:hypercube}

\begin{definition}[Boolean hypercube]\label{def:hypercube}
The \emph{Boolean hypercube}\index{hypercube} of dimension $n$ is the set
\[
  \Omega_n = \pmone^n = \{(x_1, \ldots, x_n) : x_i \in \{-1, +1\}\}.
\]
It has $|\Omega_n| = 2^n$ elements. We equip it with the \emph{uniform
probability measure}: each $x \in \Omega_n$ has probability $2^{-n}$.
\end{definition}

\begin{remark}[Why $\pmone$ and not $\{0,1\}$?]\label{rem:pm1-vs-01}
Both conventions appear in the literature. The $\pmone$ convention is
standard in Boolean Fourier analysis because it makes the algebra cleaner:
\begin{itemize}
  \item Multiplication is the natural group operation ($x_i \cdot x_j \in \pmone$),
        whereas $\{0,1\}$ would require XOR (addition mod~2).
  \item The Fourier basis functions (parity functions, below) take values
        in $\pmone$, matching the ``alphabet'' of Boolean functions.
  \item The connection to the Walsh--Hadamard transform (\Cref{ch:wht})
        is immediate.
\end{itemize}
To convert: if $b \in \{0,1\}$, set $x = 1 - 2b = (-1)^b$.
\end{remark}

\begin{definition}[Boolean and pseudo-Boolean functions]\label{def:boolean-fn}
A \emph{Boolean function}\index{Boolean function} is a map $f \colon \pmone^n \to \pmone$.
A \emph{pseudo-Boolean function}\index{pseudo-Boolean function} is a map $f \colon \pmone^n \to \R$.
\end{definition}

The set of all pseudo-Boolean functions $f \colon \pmone^n \to \R$ forms
a real vector space of dimension $2^n$ (one degree of freedom per point
in $\Omega_n$). We equip it with the \emph{expectation inner product}:

\begin{definition}[Inner product on functions]\label{def:fn-inner-product}
For $f, g \colon \pmone^n \to \R$:
\begin{equation}\label{eq:fn-ip}
  \ip{f}{g} = \E_x[f(x)\, g(x)]
  = \frac{1}{2^n} \sum_{x \in \pmone^n} f(x)\, g(x).
\end{equation}
The associated norm is $\norm{f}_2 = \sqrt{\ip{f}{f}}
= \sqrt{\E_x[f(x)^2]}$.
\end{definition}

\begin{example}\label{ex:boolean-norm}
If $f \colon \pmone^n \to \pmone$ is Boolean, then $f(x)^2 = 1$ for all
$x$, so $\norm{f}_2^2 = 1$. Every Boolean function is a unit vector in
this inner product space.
\end{example}


% ──────────────────────────────────────────────────────────────────────
\section{The Parity Functions}\label{sec:parity}

We now construct an orthonormal basis for the space of pseudo-Boolean
functions. The key objects are the \emph{parity functions} (also called
\emph{Walsh functions} \citep{walsh1923closed} or \emph{characters};
see \Cref{sec:characters} for the group-theoretic viewpoint).

\begin{definition}[Parity function]\label{def:parity}\index{parity function}
For each subset $S \subseteq [n] = \{1, 2, \ldots, n\}$, define
\begin{equation}\label{eq:chi-def}
  \chi_S(x) = \prod_{i \in S} x_i.
\end{equation}
By convention, $\chi_\emptyset(x) = 1$ for all $x$ (the empty product).
\end{definition}

Since each $x_i \in \pmone$, each $\chi_S(x)$ is a product of $\pm 1$
values, hence $\chi_S(x) \in \pmone$. The parity functions are themselves
Boolean functions. This is a distinctive feature of the Boolean Fourier
basis: the basis elements live in the same alphabet as the functions
being analyzed.

\begin{remark}[Analogy with classical Fourier analysis]\label{rem:fourier-analogy}
In classical Fourier analysis, the functions $\{e^{2\pi i k x}\}_{k \in \Z}$
form an orthonormal basis for $L^2([0,1])$, and every square-integrable
function decomposes into ``frequencies'' indexed by integers $k$. Here,
the parity functions $\{\chi_S\}_{S \subseteq [n]}$ play the same role:
every function on $\pmone^n$ decomposes into components indexed by subsets
$S$. (For more on this analogy, see \citealp[\S1.1]{odonnell2014analysis}.) The \emph{degree} $|S|$ is the analogue of frequency: degree-0 is
the DC component (the mean), degree-1 terms vary linearly, and high-degree
terms oscillate rapidly across the hypercube. The key simplification is
that the Boolean basis is real-valued ($\pm 1$), so there are no complex
numbers.
\end{remark}

\begin{example}\label{ex:parity-functions}
For $n = 3$:
\begin{align*}
  \chi_\emptyset(x) &= 1, \\
  \chi_{\{1\}}(x) &= x_1, \quad
  \chi_{\{2\}}(x) = x_2, \quad
  \chi_{\{3\}}(x) = x_3, \\
  \chi_{\{1,2\}}(x) &= x_1 x_2, \quad
  \chi_{\{1,3\}}(x) = x_1 x_3, \quad
  \chi_{\{2,3\}}(x) = x_2 x_3, \\
  \chi_{\{1,2,3\}}(x) &= x_1 x_2 x_3.
\end{align*}
There are $2^3 = 8$ such functions, one for each subset of $[3]$.
\end{example}


% ──────────────────────────────────────────────────────────────────────
\section{Orthonormality}\label{sec:orthonormality}

\begin{theorem}[Orthonormality of parity functions]\label{thm:orthonormality}\index{orthonormality}
For any $S, T \subseteq [n]$:
\begin{equation}\label{eq:orthonormality}
  \ip{\chi_S}{\chi_T} = \delta_{ST} =
  \begin{cases}
    1 & \text{if } S = T, \\
    0 & \text{if } S \neq T.
  \end{cases}
\end{equation}
\end{theorem}

\begin{proof}
We compute:
\[
  \ip{\chi_S}{\chi_T}
  = \E_x\!\left[\prod_{i \in S} x_i \cdot \prod_{j \in T} x_j\right]
  = \E_x\!\left[\prod_{i \in S \triangle T} x_i\right],
\]
where $S \triangle T = (S \setminus T) \cup (T \setminus S)$ is the
symmetric difference (we used $x_i^2 = 1$ to cancel shared indices).

If $S = T$, then $S \triangle T = \emptyset$, the product is~1, and the
expectation is~1.

If $S \neq T$, then $S \triangle T \neq \emptyset$; pick any
$k \in S \triangle T$. Since $x_k$ is independent of $\{x_i : i \neq k\}$
and takes $\pm 1$ with equal probability:
\[
  \E_x\!\left[\prod_{i \in S \triangle T} x_i\right]
  = \E_{x_k}[x_k] \cdot \E_{x_{-k}}\!\left[\prod_{\substack{i \in S \triangle T \\ i \neq k}} x_i\right]
  = 0 \cdot (\text{something}) = 0.
\]
(Here $x_{-k}$ denotes the coordinates other than $k$.)
\end{proof}

Since there are $2^n$ parity functions and the space of pseudo-Boolean
functions has dimension $2^n$, orthonormality implies:

\begin{corollary}\label{cor:parity-onb}
$\{\chi_S\}_{S \subseteq [n]}$ is an orthonormal basis for the space
of all functions $f \colon \pmone^n \to \R$ under the inner product
$\ip{f}{g} = \E_x[f(x)g(x)]$.
\end{corollary}


% ──────────────────────────────────────────────────────────────────────
\section{The Fourier Expansion}\label{sec:fourier-expansion}

By \Cref{prop:fourier-finite} (general Fourier expansion) and
\Cref{cor:parity-onb} (the parity functions form an ONB), we immediately
obtain:

\begin{theorem}[Fourier expansion on $\pmone^n$]\label{thm:fourier-expansion}\index{Fourier expansion}
Every function $f \colon \pmone^n \to \R$ has a unique expansion
\begin{equation}\label{eq:fourier-expansion}
  f(x) = \sum_{S \subseteq [n]} \fhat(S)\, \chi_S(x),
\end{equation}
where the \emph{Fourier coefficients}\index{Fourier coefficient} (or \emph{Walsh--Fourier
coefficients}) are
\begin{equation}\label{eq:fourier-coeff}
  \fhat(S) = \ip{f}{\chi_S}
  = \E_x\!\left[f(x) \prod_{i \in S} x_i\right]
  = \frac{1}{2^n} \sum_{x \in \pmone^n} f(x) \prod_{i \in S} x_i.
\end{equation}
\end{theorem}

The expansion \eqref{eq:fourier-expansion} is a \emph{multilinear
polynomial}: each variable $x_i$ appears with degree at most~1 (because
$x_i^2 = 1$ for $x_i \in \pmone$, so any higher power reduces). This is
the unique multilinear representation of $f$.

\begin{definition}[Degree]\label{def:degree}
The \emph{degree}\index{degree (Fourier)} of the monomial $\chi_S$ is $|S|$. We say $f$ has
\emph{Fourier degree at most $d$} if $\fhat(S) = 0$ for all $S$ with
$|S| > d$.
\end{definition}


% ──────────────────────────────────────────────────────────────────────
\section{Parseval's Identity}\label{sec:parseval}

Applying \Cref{prop:fourier-finite} (Parseval in the abstract setting)
with the parity ONB:

\begin{theorem}[Parseval's identity]\label{thm:parseval}\index{Parseval's identity}
For any $f \colon \pmone^n \to \R$:
\begin{equation}\label{eq:parseval}
  \norm{f}_2^2 = \E_x[f(x)^2] = \sum_{S \subseteq [n]} \fhat(S)^2.
\end{equation}
\end{theorem}

For Boolean functions ($f \colon \pmone^n \to \pmone$):
\begin{equation}\label{eq:parseval-boolean}
  \sum_{S \subseteq [n]} \fhat(S)^2 = 1.
\end{equation}

This is the \emph{energy constraint}: the Fourier coefficients of a
Boolean function lie on the unit sphere in $\R^{2^n}$. The ``spectral
energy'' $\fhat(S)^2$ at each subset $S$ can be thought of as the
fraction of the function's total variance explained by that frequency
component.

\begin{remark}[Relevance to the thesis]\label{rem:parseval-thesis}
The thesis uses Parseval's identity in a crucial way: it implies that the
Fourier coefficients of any Boolean function are automatically normalized
($\sum_S \fhat(S)^2 = 1$). This means the spectral parameterization of
binary weight matrices comes with a built-in energy budget, replacing the
learned per-layer scaling factors $\alpha$ used in BitNet.
\end{remark}


% ──────────────────────────────────────────────────────────────────────
\section{Worked Examples}\label{sec:fourier-examples}

We compute the Fourier expansion for several important Boolean functions.
These examples build intuition for ``spectral concentration''---the
phenomenon that many natural Boolean functions have most of their Fourier
mass on low-degree coefficients.

\subsection{Dictator Functions}\label{sec:dictator}

\begin{example}[Dictator]\label{ex:dictator}
The $i$-th \emph{dictator function}\index{dictator function} is $f(x) = x_i$. Its Fourier
expansion is trivially $f = \chi_{\{i\}}$, so
\[
  \fhat(S) = \begin{cases} 1 & \text{if } S = \{i\}, \\ 0 & \text{otherwise.} \end{cases}
\]
All spectral energy is at degree~1. Parseval: $1^2 = 1$. \checkmark
\end{example}

Dictator functions are the extreme case of spectral concentration: all
energy is on a single degree-1 coefficient.


\subsection{The AND Function}\label{sec:and}

\begin{example}[AND on 2 bits]\label{ex:and}\index{AND function}
Define $\text{AND}(x_1, x_2) = +1$ iff $x_1 = x_2 = +1$, and $-1$
otherwise. The truth table:
\begin{center}
\begin{tabular}{ccc}
\toprule
$x_1$ & $x_2$ & $\text{AND}$ \\
\midrule
$-1$ & $-1$ & $-1$ \\
$-1$ & $+1$ & $-1$ \\
$+1$ & $-1$ & $-1$ \\
$+1$ & $+1$ & $+1$ \\
\bottomrule
\end{tabular}
\end{center}

Computing coefficients using \eqref{eq:fourier-coeff}, $\fhat(S) = \frac{1}{4}\sum_x f(x)\,\chi_S(x)$:
\begin{align*}
  \fhat(\emptyset) &= \tfrac{1}{4}[(-1)+(-1)+(-1)+(+1)] = -\tfrac{1}{2}, \\
  \fhat(\{1\}) &= \tfrac{1}{4}[(-1)(-1) + (-1)(-1) + (-1)(+1) + (+1)(+1)]
  = \tfrac{1}{4}[1 + 1 - 1 + 1] = \tfrac{1}{2}, \\
  \fhat(\{2\}) &= \tfrac{1}{4}[(-1)(-1) + (-1)(+1) + (-1)(-1) + (+1)(+1)]
  = \tfrac{1}{4}[1 - 1 + 1 + 1] = \tfrac{1}{2}, \\
  \fhat(\{1,2\}) &= \tfrac{1}{4}[(-1)(+1) + (-1)(-1) + (-1)(-1) + (+1)(+1)]
  = \tfrac{1}{4}[-1+1+1+1] = \tfrac{1}{2}.
\end{align*}

So: $\text{AND}(x_1, x_2) = -\frac{1}{2} + \frac{1}{2}x_1 + \frac{1}{2}x_2 + \frac{1}{2}x_1 x_2$.

Parseval: $\frac{1}{4} + \frac{1}{4} + \frac{1}{4} + \frac{1}{4} = 1$. \checkmark

The spectral energy is spread evenly across degrees 0, 1, 1, 2.%
\footnote{See companion notebook, \S1: runnable verification of AND and
Maj$_3$ coefficients, Parseval checks, degree truncation, and the
convolution theorem.}
\end{example}


\subsection{The Majority Function}\label{sec:majority}

\begin{example}[Majority on 3 bits]\label{ex:majority}\index{majority function}
$\text{Maj}_3(x) = \sgn(x_1 + x_2 + x_3)$, where $\sgn(z) = +1$ if
$z > 0$ and $\sgn(z) = -1$ if $z < 0$ (for odd $n$, the sum
$x_1 + \cdots + x_n$ is never zero, so $\sgn(0)$ does not arise).
This outputs $+1$ if at least 2 of 3 inputs are $+1$.

Truth table and Fourier computation:
\begin{center}
\begin{tabular}{cccc}
\toprule
$x_1$ & $x_2$ & $x_3$ & $\text{Maj}_3$ \\
\midrule
$-1$ & $-1$ & $-1$ & $-1$ \\
$-1$ & $-1$ & $+1$ & $-1$ \\
$-1$ & $+1$ & $-1$ & $-1$ \\
$-1$ & $+1$ & $+1$ & $+1$ \\
$+1$ & $-1$ & $-1$ & $-1$ \\
$+1$ & $-1$ & $+1$ & $+1$ \\
$+1$ & $+1$ & $-1$ & $+1$ \\
$+1$ & $+1$ & $+1$ & $+1$ \\
\bottomrule
\end{tabular}
\end{center}

The Fourier coefficients are computed by the same procedure as the AND
example (applying \eqref{eq:fourier-coeff} to each of the 8 rows):
\begin{align*}
  \fhat(\emptyset) &= 0 \quad \text{(balanced: 4 outputs $+1$, 4 outputs $-1$)}, \\
  \fhat(\{1\}) = \fhat(\{2\}) = \fhat(\{3\}) &= \tfrac{1}{2}, \\
  \fhat(\{1,2\}) = \fhat(\{1,3\}) = \fhat(\{2,3\}) &= 0, \\
  \fhat(\{1,2,3\}) &= -\tfrac{1}{2}.
\end{align*}

Hence: $\text{Maj}_3(x) = \frac{1}{2}(x_1 + x_2 + x_3) - \frac{1}{2}x_1 x_2 x_3$.

Parseval: $3 \cdot \frac{1}{4} + \frac{1}{4} = 1$. \checkmark

The spectral energy breakdown by degree:
\begin{center}
\begin{tabular}{lcc}
\toprule
Degree & Energy & Fraction \\
\midrule
0 & 0 & 0\% \\
1 & $3/4$ & 75\% \\
2 & 0 & 0\% \\
3 & $1/4$ & 25\% \\
\bottomrule
\end{tabular}
\end{center}

$\text{Maj}_3$ is \emph{spectrally concentrated}: 75\% of its energy is
at degree~1. This is a general phenomenon: the majority function on $n$
bits has most of its spectral energy on degree-1 coefficients, with the
concentration improving as $n$ grows
\citep[Chapter~5]{odonnell2014analysis}.
\end{example}


\subsection{The Parity Function (Anti-Concentration)}\label{sec:parity-ex}

\begin{example}[Parity]\label{ex:parity}
$f(x) = x_1 x_2 \cdots x_n = \chi_{[n]}(x)$.

This function has $\fhat([n]) = 1$ and all other coefficients zero. All
energy is at the maximum degree $n$. The parity function is the
\emph{opposite} of spectrally concentrated: it is maximally complex in the
Fourier sense. It is also the hardest Boolean function to learn from
random examples \citep{linial1993constant}.
\end{example}


% ──────────────────────────────────────────────────────────────────────
\section{Degree-Level Energy Decomposition}\label{sec:energy-levels}

Parseval allows us to decompose the ``energy'' of a Boolean function by
degree:

\begin{definition}[Spectral energy at degree $k$]\label{def:energy-level}
For $f \colon \pmone^n \to \R$ and $0 \leq k \leq n$:
\[
  W^k[f] = \sum_{\substack{S \subseteq [n] \\ |S| = k}} \fhat(S)^2.
\]
The \emph{spectral energy profile} is the vector $(W^0[f], W^1[f], \ldots, W^n[f])$.
\end{definition}

By Parseval: $\sum_{k=0}^n W^k[f] = \norm{f}_2^2$. For Boolean $f$,
this sum is~1, so the energy profile is a probability distribution on
degrees $\{0, 1, \ldots, n\}$.

For a \emph{random} Boolean function (chosen uniformly at random from all
$2^{2^n}$ Boolean functions on $n$ bits), the expected energy at degree
$k$ is $\binom{n}{k} / 2^n$.  To see this: for a random Boolean $f$,
$\E[\fhat(S)^2] = \E[\E_x[f(x)\chi_S(x)]^2]$. By expanding the square
and using independence of $f(x)$ and $f(x')$ for $x \neq x'$, one
obtains $\E[\fhat(S)^2] = 1/2^n$ for every $S \neq \emptyset$. Since
there are $\binom{n}{k}$ subsets of size $k$, the expected energy at
degree $k$ is $\binom{n}{k}/2^n$, which follows a
$\text{Binomial}(n, 1/2)$ distribution peaked at $k = n/2$.
So a random Boolean function has its energy concentrated around degree
$n/2$. (See \citealp[Exercise~1.19]{odonnell2014analysis} for this
calculation.)

The thesis's central hypothesis is that \emph{trained} binary weight
matrices, viewed as Boolean functions, have their energy concentrated at
much lower degrees---far below the $n/2$ baseline of random functions.


% ──────────────────────────────────────────────────────────────────────
\section{Low-Degree Truncation}\label{sec:truncation}

\begin{definition}[Degree-$k$ truncation]\label{def:truncation}
For $f \colon \pmone^n \to \R$:
\[
  f^{\leq k}(x) = \sum_{\substack{S \subseteq [n] \\ |S| \leq k}} \fhat(S)\, \chi_S(x).
\]
Similarly, $f^{>k} = f - f^{\leq k}$.
\end{definition}

The truncation $f^{\leq k}$ is the best degree-$k$ approximation to $f$
in the $L^2$ norm---this is a standard consequence of orthogonal
projection in Hilbert spaces (see \citealp[Chapter~6]{axler2024linear}):
the ``tail'' $f^{>k}$ is orthogonal to all degree-$\leq k$ functions. The
approximation error is:
\begin{equation}\label{eq:truncation-error}
  \norm{f - f^{\leq k}}_2^2 = \norm{f^{>k}}_2^2
  = \sum_{|S| > k} \fhat(S)^2.
\end{equation}

For Boolean functions, if we want to \emph{round} $f^{\leq k}$ back to
$\pmone$ via the sign function, the fraction of inputs where the rounded
version disagrees with $f$ is bounded by:

\begin{proposition}[Sign-rounding error; thesis Proposition~3.1]\label{prop:sign-rounding}\index{sign-rounding error}
Let $f \colon \pmone^n \to \pmone$ be Boolean and
$\tilde{f} = \sgn(f^{\leq k})$. Then
\[
  \Prob_x[\tilde{f}(x) \neq f(x)] \leq \sum_{|S| > k} \fhat(S)^2.
\]
\end{proposition}

\begin{proof}
If $\tilde{f}(x) \neq f(x)$, then $f(x) \cdot f^{\leq k}(x) \leq 0$.
Write $f^{\leq k}(x) = f(x) - f^{>k}(x)$, so
$f(x) \cdot f^{\leq k}(x) = 1 - f(x) \cdot f^{>k}(x)$.
Thus the disagreement requires $f(x) \cdot f^{>k}(x) \geq 1$, which
implies $|f^{>k}(x)| \geq 1$. By Markov's inequality
(\Cref{thm:markov}) applied to $(f^{>k})^2$:
\[
  \Prob[|f^{>k}(x)| \geq 1] \leq \E[(f^{>k}(x))^2] = \sum_{|S|>k} \fhat(S)^2.
  \qedhere
\]
\end{proof}

This bound is tight when $f^{>k}$ is small everywhere it is nonzero
(as happens for spectrally concentrated functions), and loose when
$f^{>k}$ takes large values on a small set.

\begin{example}[Truncation of $\text{Maj}_3$]\label{ex:maj3-truncation}
For $\text{Maj}_3$, the tail energy above degree 1 is
$\fhat(\{1,2,3\})^2 = 1/4$, so the bound gives
$\Prob[\sgn(f^{\leq 1}) \neq f] \leq 1/4$. In fact, $f^{\leq 1}(x) =
\frac{1}{2}(x_1 + x_2 + x_3)$, which has the same sign as $\text{Maj}_3$
on all 8 inputs (since $|x_1 + x_2 + x_3| \geq 1$ always), so the
actual error is $0/8 = 0$. The bound is not tight here because
$f^{>1}(x) = -\frac{1}{2}x_1 x_2 x_3$ never has magnitude exceeding
$\frac{1}{2} < 1$.
\end{example}


% ──────────────────────────────────────────────────────────────────────
\section{Plancherel and the Convolution Theorem}\label{sec:convolution}

Two additional Fourier identities that are used in the thesis:

\begin{theorem}[Plancherel's identity]\label{thm:plancherel}
For $f, g \colon \pmone^n \to \R$:
\[
  \ip{f}{g} = \sum_{S \subseteq [n]} \fhat(S)\, \ghat(S).
\]
\end{theorem}

\begin{proof}
Expand both functions:
$\ip{f}{g} = \E_x\!\left[\sum_S \fhat(S)\chi_S(x) \cdot \sum_T \ghat(T)\chi_T(x)\right]
= \sum_S \sum_T \fhat(S)\,\ghat(T)\,\E_x[\chi_S(x)\chi_T(x)]
= \sum_S \fhat(S)\,\ghat(S)$,
where the last step uses orthonormality: $\E[\chi_S \chi_T] = \delta_{ST}$.
\end{proof}

\begin{definition}[Dyadic convolution]\label{def:convolution}
For $f, g \colon \pmone^n \to \R$, the \emph{dyadic convolution}\index{convolution!dyadic} is
\[
  (f \ast g)(x) = \E_y[f(y)\, g(x \cdot y)],
\]
where $x \cdot y = (x_1 y_1, \ldots, x_n y_n)$ is coordinatewise
multiplication. (The name ``dyadic'' refers to the group $\Z_2^n$;
see \Cref{sec:group-theory}.)
\end{definition}

\begin{theorem}[Convolution theorem]\label{thm:convolution}
$\widehat{f \ast g}(S) = \fhat(S) \cdot \ghat(S)$.
\end{theorem}

\begin{proof}
\begin{align*}
  \widehat{f \ast g}(S)
  &= \E_x[(f \ast g)(x)\, \chi_S(x)]
  = \E_x\!\left[\E_y[f(y)\, g(x \cdot y)]\, \chi_S(x)\right] \\
  &= \E_y\!\left[f(y)\, \E_x[g(x \cdot y)\, \chi_S(x)]\right].
\end{align*}
Substituting $z = x \cdot y$ (which is uniform when $x$ is uniform,
since $y$ is fixed and coordinatewise multiplication by $y \in \pmone^n$
is a bijection on $\pmone^n$):
\[
  \E_x[g(x \cdot y)\, \chi_S(x)]
  = \E_z[g(z)\, \chi_S(z \cdot y)]
  = \E_z[g(z)\, \chi_S(z)\, \chi_S(y)]
  = \ghat(S)\, \chi_S(y),
\]
using $\chi_S(z \cdot y) = \chi_S(z)\,\chi_S(y)$ (since $\chi_S$ is
multiplicative: $\prod_{i \in S}(z_i y_i) = \prod_{i \in S} z_i \cdot \prod_{i \in S} y_i$).

Therefore:
$\widehat{f \ast g}(S) = \ghat(S)\, \E_y[f(y)\, \chi_S(y)] = \ghat(S)\, \fhat(S)$.
\end{proof}


% ──────────────────────────────────────────────────────────────────────
\section*{Recommended Reading}

\begin{itemize}
  \item \textbf{\citet{odonnell2014analysis}, \emph{Analysis of Boolean
    Functions}.}
    The definitive textbook. Chapters~1--2 cover everything in this
    chapter plus exercises and historical notes. Free preprint on arXiv
    (2105.10386). Graduate level; assumes real analysis and basic
    probability.

  \item \textbf{\citet{dewolf2008brief}, ``A Brief Introduction to
    Fourier Analysis on the Boolean Cube.''}
    A 30-page survey covering the core definitions, Parseval, influences,
    and a few applications. Excellent for a first pass before reading
    O'Donnell. Freely available online. Assumes undergraduate
    linear algebra.

  \item \textbf{\citet{mansour1994learning}, ``Learning Boolean Functions
    via the Fourier Transform.''}
    An early survey emphasizing the learning-theory applications of
    Boolean Fourier analysis. Good for seeing how Fourier coefficients
    connect to PAC learning. Assumes familiarity with basic
    computational complexity.
\end{itemize}
